\section*{Capítulo \ref{ch:b3:asymmetric-synthesis} --- Síntese assimétrica}

\begin{solution}{exo:b3:asymmetric-synthesis:1}
ee 80\,\%: $\mathrm{er} = 1.8/0.2 = 9$ (90\,:\,10). er 99\,:\,1: $\mathrm{ee} = 98/100 = 98$\,\%. ee 50\,\%: 75\,\% e 25\,\%.
\end{solution}

\begin{solution}{exo:b3:asymmetric-synthesis:2}
$(1840 - 60)/(1840 + 60) = 0.937$: ee 93.7\,\%, um er de $1840/60 = 31$.
\end{solution}

\begin{solution}{exo:b3:asymmetric-synthesis:3}
Prioridades O $>$ etila $>$ H. Desenhado com O para cima, etila embaixo à esquerda e H embaixo à direita, a
face da frente é \textit{Si}, a de trás, \textit{Re}. Um grupo metila adicionado pela \omterm{def:b3:asymmetric-synthesis:faces}{face Re}, com
as prioridades OH $>$ etila $>$ metila $>$ H no produto, dá o (\textit{R})-butan-2-ol.
\end{solution}

\begin{solution}{exo:b3:asymmetric-synthesis:4}
Etanol: \omterm{def:b3:asymmetric-synthesis:topicity}{enantiotópicos} (trocados pelo \omterm{def:b3:point-groups:mirror}{plano especular} da molécula; mesmo
deslocamento em um solvente aquiral). Butan-2-ol: \omterm{def:b3:asymmetric-synthesis:topicity}{diastereotópicos} (vizinhos a um estereocentro;
deslocamentos diferentes).
\end{solution}

\begin{solution}{exo:b3:asymmetric-synthesis:5}
er 199: $RT\ln 199 = \qty{13.1}{kJ/mol}$ a \qty{298}{K}, \qty{8.6}{kJ/mol} a \qty{195}{K}. 80\,\% de ee a \qty{25}{\celsius} é um er de 9,
$\Delta\Delta G^\ddagger = RT\ln 9 = \qty{5.45}{kJ/mol}$; a \qty{-78}{\celsius}, $\mathrm{er} = \eu^{5450/(8.314 \times 195.15)} = 28.7$, ee 93\,\%.
\end{solution}

\begin{solution}{exo:b3:asymmetric-synthesis:6}
\omterm{def:b3:asymmetric-synthesis:ee}{Pureza óptica} $12.0/15.0 = 80$\,\%: 90\,\% de (\textit{S}) e 10\,\% de (\textit{R}). Um poder rotatório pequeno, as impurezas,
os efeitos de concentração e de solvente e um comportamento não linear entre poder rotatório e composição
a tornam pouco confiável; a cromatografia separa e conta os
enantiômeros diretamente.
\end{solution}

\begin{solution}{exo:b3:asymmetric-synthesis:7}
$s = \ln(0.40 \times 0.10)/\ln(0.40 \times 1.90) = -3.22/-0.274 = 12$.
\end{solution}

\begin{solution}{exo:b3:asymmetric-synthesis:8}
Com o L-($+$)-tartarato de dietila, (2\textit{S},3\textit{S})-2,3-epoxi-hexan-1-ol; com o D-($-$)-tartarato de dietila, o
enantiômero, (2\textit{R},3\textit{R}).
\end{solution}

\begin{solution}{exo:b3:asymmetric-synthesis:9}
L = fenila, M = metila, S = H. Para o aldeído (\textit{S}), o confôrmero reativo tem a fenila
perpendicular ao plano da carbonila e o hidrogênio ao lado do caminho do
nucleófilo, que ataca em anti à fenila. A determinação das configurações dá
o (2\textit{S},3\textit{S})-3-fenilbutan-2-ol como diastereoisômero majoritário.
\end{solution}

\begin{solution}{exo:b3:asymmetric-synthesis:10}
No início, os dois enantiômeros estão presentes em quantidades iguais, de modo que os produtos
se formam na razão $k_{\mathrm{fast}} : k_{\mathrm{slow}} = s$: $\mathrm{ee}_p = (s - 1)/(s + 1)$, 0.905 para $s = 20$. À medida que a reação avança, o
enantiômero rápido se esgota e o ee do produto cai, enquanto o ee do substrato sobe
em direção a 1: o substrato sempre pode ser purificado indo mais longe; o produto
não.
\end{solution}

\begin{solution}{exo:b3:asymmetric-synthesis:11}
Com racemização rápida: até 100\,\% de rendimento, $\mathrm{er} = 99$, ee 98\,\%. Sem ela, a 50\,\% de
conversão: rendimento de no máximo 50\,\% e um ee do produto de 93\,\% (o ee do substrato é o
mesmo, 93\,\%).
\end{solution}

\begin{solution}{exo:b3:asymmetric-synthesis:12}
Quando os dois complexos se interconvertem mais depressa do que reagem, a razão dos produtos é
$k_B[\mathrm B]/(k_A[\mathrm A]) = 600/10 = 60$: o complexo minoritário B dá 98\,\% do produto (ee 97\,\%). A
seletividade é fixada pela diferença entre os dois estados de transição, e não pelas
\omterm{def:b3:partition-functions:boltzmann}{populações} dos intermediários.
\end{solution}

\begin{solution}{pb:b3:asymmetric-synthesis:1}
\textbf{1.} O carbono do alqueno que porta o nitrogênio tem três substituintes diferentes
e duas faces diferentes; o hidrogênio adicionado a ele faz do carbono $\alpha$ do
aminoácido um estereocentro.

\textbf{2.} O racemato: as duas faces são atacadas com a mesma velocidade.

\textbf{3.} Seu oxigênio carbonílico se liga ao ródio junto com a ligação C=C: o substrato é
mantido como quelato, em uma geometria fixa que o ligante pode discriminar.

\textbf{4.} Ela constrói um bolsão quiral em torno do metal (de simetria $C_2$): as duas maneiras de
ligar o substrato, por uma face ou pela outra, são diastereoisoméricas, com energias e
reatividades diferentes.

\textbf{5.} Não: o complexo minoritário reage muito mais depressa com o hidrogênio (Curtin--Hammett).

\textbf{6.} O metal fica em uma das faces do alqueno ligado, e os hidrogênios são
transferidos a partir do metal.

\textbf{7.} $\mathrm{er} = 1.95/0.05 = 39$.

\textbf{8.} $\Delta\Delta G^\ddagger = RT\ln 39 = 8.314 \times 298.15 \times 3.664 = \qty{9.1}{kJ/mol}$.

\textbf{9.} $\mathrm{er} = \eu^{9082/(8.314 \times 195.15)} = 270$: ee 99.3\,\%.

\textbf{10.} $\mathrm{er} = \eu^{11082/(8.314 \times 298.15)} = 87$: ee 97.7\,\%.

\textbf{11.} Menos que uma ligação de hidrogênio típica: uma pequena diferença nos contatos
estéricos ou uma única interação fraca entre substrato e ligante decide o
resultado, e é por isso que esses catalisadores são difíceis de projetar a partir de primeiros princípios.

\textbf{12.} A velocidade cai, o hidrogênio se dissolve e o catalisador se comporta de outro modo
a baixa temperatura, e resfriar um grande reator custa energia.

\textbf{13.} 97.5 g do enantiômero (\textit{S}) e 2.5 g do (\textit{R}).

\textbf{14.} $97.5 - 5.0 = \qty{92.5}{g}$ de cristais, praticamente enantiopuros.

\textbf{15.} $92.5/100 = 92.5$\,\% do produto bruto (95\,\% do enantiômero (\textit{S})).

\textbf{16.} Os cristais recolhem o enantiômero em excesso, e a parte racêmica permanece em
solução com o minoritário; a partir de um racemato não há excesso a recolher, e a
cristalização dá cristais racêmicos.

\textbf{17.} Por HPLC quiral, com uma referência racêmica.

\textbf{18.} 50\,\%: só o enantiômero lento sobrevive.

\textbf{19.} $s = \ln(0.45 \times 0.05)/\ln(0.45 \times 1.95) = -3.79/-0.131 = 29$.

\textbf{20.} 45\,\% do racemato (88\,\% do enantiômero desejado presente no início,
com 95\,\% de ee).

\textbf{21.} Com racemização rápida do substrato, o rendimento poderia se aproximar de 100\,\%.

\textbf{22.} Ela converte todo o precursor aquiral barato no enantiômero correto,
com um catalisador usado em quantidade ínfima e sem descartar nenhum enantiômero.

\textbf{23.} O ródio é tóxico, e seus resíduos nos fármacos são limitados a algumas partes por
milhão; eles são medidos por espectrometria de massa com plasma indutivamente acoplado.

\textbf{24.} 95\,\% de ee a \qty{25}{\celsius} correspondem a $\Delta\Delta G^\ddagger = RT\ln 39 \approx \qty{9.1}{kJ/mol}$.
\end{solution}
